\lesson{12/10}
\subsection{Charged scalar field}
For the simple scalar field considered earlier, no internal symmetries can be
found (?): to generalize the concept, one must consider a complex field, which has
multiple components while retaining the fact that it is scalar:
consider two independent real scalar fields $\phi_1,\phi_2$ and their linear
combinations
\begin{gather*}
\phi(x^\mu) = \frac{1}{\sqrt{2}}\left(\phi_1+i\phi_2\right)\\
\phi^\dagger(x^\mu) = \frac{1}{\sqrt{2}}\left(\phi_1-i\phi_2\right)
\end{gather*}
The simplest Lagrangian using this kind of field presents the simplest bilinear
Lorentz invariance:
$$\mathcal{L}=\pd_\mu\phi\pd^\mu\phi^\dagger - m^2\phi^\dagger\phi$$
The Euler-Lagrange equations turn out to be
$$\begin{dcases*}
  \left(\Box + m^2\right)\phi(x^\mu) = 0\\
  \left(\Box + m^2\right)\phi^\dagger(x^\mu) = 0
\end{dcases*}.$$
Consider the internal transformation that gives an arbitrary phase shift to the field
\begin{gather*}
\phi \rightarrow \phi'=e^{i\alpha}\phi\rightarrow \phi'=\phi+i\alpha\phi\\
\phi\rightarrow{\phi^\dagger}'=e^{i\alpha}\phi\rightarrow{\phi^\dagger}'=\phi^\dagger
+ i\alpha\phi^\dagger
\end{gather*}
this is a symmetry with respect to the group $U(1)$.

The Noether current is
$$J^\mu = \pdv[]{\mathcal{L}}{\pd_\mu\phi}\delta\phi + \pdv{\mathcal{L}}
{\pd_\mu\phi^\dagger}\delta\phi^\dagger,$$
expanding the terms and factorizing the parameter $\alpha$ yields
$$J^\mu = \pd^\mu\phi^\dagger(i\phi)+\pd^\mu\phi(-i\phi^\dagger) = 
i\left(\phi^\dagger\pd_\mu\phi - \phi\pd_\mu\phi^\dagger\right).$$
Noether's theorem states that the quantity
$$Q = \int_{}^{}\dd^3 x J^0 = i\int_{}^{}\dd^3x \left(\phi^\dagger\dot\phi-
\phi\dot\phi^\dagger\right)$$
is conserved: in classical field theory, this does not seem like much, but with
canonical quantization it will be clear that this $Q$ is \textbf{the total charge of
the system}.
\section{Spinor Field}
In this section the spinor field will be defined and its invariances will be studied.
\subsection{Spinor representations of the Lorentz Group}
As hinted before, there are 2 irreducible representations associated with spinor
fields, both related to the $\mathfrak{su}(2)$ algebra:
\begin{gather*}
    \left(\frac{1}{2},0\right) \rightarrow \psi_L(x^\mu) =
    \begin{pmatrix}
    \psi_L^1(x^\mu)\\
    \psi_L^2(x\mu)
    \end{pmatrix}
    \\
    \phantom{x}\\
    \left(0,\frac{1}{2}\right) \rightarrow \psi_R(x^\mu)=
    \begin{pmatrix}
    \psi_L^1(x^\mu)\\
    \psi_L^2(x^\mu)
    \end{pmatrix}\\
\end{gather*}
Consider now the Lorentz transformations $\Lambda_L$ and $\Lambda_R$, which are
$2\times2$ matrices and are expected to be generated by the angular momenta $J^i$
and boosts $K^i$ operators, with the caveat, given by Wigner's theorem, that,
since they do not form a compact group, the $K$'s do not admit unitary
representation:
\begin{gather*}
J^i = \frac{1}{2}\sigma^i\rightarrow \exp\left[\frac{i}{2}\omega_i\sigma^i\right]\\
K^i = \mp i \frac{\sigma^i}{2}\rightarrow \exp\left[\mp\frac{1}{2}\eta_i\sigma^i
\right]
\end{gather*}
Where the $\mp$ sign distinguishes
$$K_L^i = -i\frac{\sigma^i}{2},\,K_R^i = +i\frac{\sigma^i}{2}.$$
So the Lorentz transformations are 
\begin{gather*}
\Lambda_L = \exp\left[\frac{1}{2}\sigma^i\left(i\omega_i-\eta_i\right)\right]\\
\Lambda_R = \exp\left[\frac{1}{2}\sigma^i\left(i\omega_i+\eta_i\right)\right].
\end{gather*}
These transformation are related:
$$\Lambda_L^{-1} = \Lambda_R^\dagger,\; \sigma^2\Lambda_L\sigma^2,\;
\Lambda_L^T\sigma^2\Lambda_L =\sigma^2,\;\Lambda_L^T=\sigma^2 \Lambda_L^{-1}\sigma_2$$
and the relations hold switching $R$ and $L$; these properties produce an interesting
result
$$\sigma^2\psi^*_L\underset{\text{L.T.}}\rightarrow\sigma^2\Lambda_L^*\psi^*_L = 
\sigma^2\Lambda_L^*\sigma^2\sigma^2\psi_L^* = \Lambda_R\sigma^2\psi_L^*$$
which means that $\sigma^2\psi^*_L$ transforms like an element of the
$\left(0,\frac{1}{2}\right)$ representation.
\subsection{Invariants built from spinors}
From QM1 spin composition, it is known that
$$\left(\frac{1}{2},0\right)\otimes\left(0,\frac{1}{2}\right) = (0,0)\oplus(1,0):$$
consider two left spinors $\psi_L,\chi_L$ and the object
$$\chi^T\sigma^2\psi_L,$$
how does it transform?
$$\chi^T\sigma^2\psi_L\underset{\text{L.T.}}\rightarrow\chi_L^T\Lambda_L^T\sigma^2
\Lambda_L\psi_L = \chi^T\sigma^2\psi_L$$
This combination of two spinors of the left representation is a scalar; consider the
special case $\psi_L =\chi_L$:
$$\psi_L^T\sigma^2\psi_L = \begin{pmatrix}
    \psi_L^1&\psi_L^2
\end{pmatrix}\begin{pmatrix}
    0&-i\\
    i&0
\end{pmatrix}
\begin{pmatrix}
\psi_L^1\\
\psi_L^2
\end{pmatrix}
= -i \psi_L^1\psi_L^2 + i \psi_L^2\psi_L^1
$$
in order for this expression to be non-zero, the components of the spinors must be
anticommuting numbers, called \textit{Grassman Numbers}.

Considering $\chi_L = \sigma^2\psi_R^*$ and multiplying by $i$, yields
$$i(\sigma^2\psi_R^*)^T\sigma^2\psi_L = -i\psi_R^\dagger\psi_L$$
Equivalently, one can prove that another invariant is
$$\psi_L^\dagger\psi_R.$$\\

Another way to determine invariants uses the 4D extension of the Pauli matrices
\begin{gather*}
\sigma^\mu = (\mathbb{I},\sigma^i), \\
\bar\sigma^\mu = (\mathbb{I}, -\sigma^i)
\end{gather*}
to build 4-vectors like
\begin{gather*}
\psi_L^\dagger\bar\sigma^\mu\psi_L=(\psi_L^\dagger\psi_L,
-\psi_L^\dagger\sigma^i\psi_L)\\
\psi_R^\dagger\sigma^\mu\psi_R=(\psi_R^\dagger\psi_R,
\psi_R^\dagger\sigma^i\psi_R):
\end{gather*}
one can use derivatives contracted with the index to create invariants:
\begin{gather*}
\psi_L^\dagger\bar\sigma^\mu\pd_\mu\psi_L+\psi^\dagger_R\sigma^\mu\pd_\mu\psi_R
\end{gather*}

\textbf{Note:} All these invariant are linear combinations of the two
representations: this hints at
the fact that they should not be treated separately.
\subsection{Dirac Spinor}
There is a problem with the 2 spinor representations, since the $J$'s are
pseudo-vectors and the $K$'s are not, a parity transformation acts differently on
the generators and on the
Casimir operators: it interchanges the operators $N^2, N^{\dagger2}$ and so the
eigenvalues that label the representations
$$(n,m)\underset{\text{Parity}}\longleftrightarrow(m,n).$$
This means that one of these representations is not enough to describe a spin 1/2
field which is invariant under parity transformations (because reality is!). 

The needed field is the \textbf{Dirac Spinor}, defined in the \textit{Weyl
representation} as
$$\psi = \begin{pmatrix}
\psi_L\\ \psi_R
\end{pmatrix} = \begin{pmatrix}
\psi^1_L \\ \psi^2_L \\ \psi^1_R \\ \psi^2_R
\end{pmatrix}
$$
\subsection{Expanding the group}
To unite the two spinors, a broader group needs to be used: recalling the
definition of the sectors of the Lorentz Group, the needed group is the set of
proper and improper orthochronous Lorentz transformations
$$L^\uparrow = L^\uparrow_+\cup L^\uparrow_-,$$
which contains parity operators.
In this group, the sum of the two representations of the restricted Lorentz group
is \textit{an irreducible representation}: the generators are the $4$ Dirac
$\gamma^\mu$ matrices, which in the Weyl representations are
\begin{gather*}
\gamma^\mu =\begin{pmatrix}
    \mathbb{O}&\sigma^\mu\\
    \bar\sigma^\mu&\mathbb{O}
\end{pmatrix}\longrightarrow
\gamma_0 = \begin{pmatrix}
    \mathbb{O}&\mathbb{I}\\
    \mathbb{I}&\mathbb{O}
\end{pmatrix},\,\gamma^i = \begin{pmatrix}
    \mathbb{O}&\sigma^i\\
    \sigma^i&\mathbb{O}
\end{pmatrix}
\end{gather*}

